% TOP_PROBLEM: {"id": 2146, "title": "Erdős Problem #501", "category_id": 2, "category": {"id": 2, "name": "combinatorics", "display_name": "Combinatorics", "description": "Counting problems, graph theory, discrete structures.", "slug": "combinatorics", "order_index": 2, "created_at": "2026-07-31T15:26:25.671Z"}, "status": "open"}
% FIRST_POSTED: null
% ENTRY_KIND: statement_only
% Imported from: unsolvedmath_additional_500.tex; UnsolvedMath ID 2146
% !TeX program = lualatex
% Compile twice: lualatex 2146.tex
% Self-contained: no external bibliography, images, or \input files.
% Numeric section labels and visible numbers are original UnsolvedMath IDs.
\documentclass[11pt,a4paper]{article}
\usepackage[margin=24mm,headheight=14pt]{geometry}
\usepackage{fontspec}
\setmainfont{Latin Modern Roman}
\setsansfont{Latin Modern Sans}
\setmonofont{Latin Modern Mono}
\usepackage{amsmath,amssymb,mathtools}
\usepackage{unicode-math}
\setmathfont{Latin Modern Math}
\usepackage{microtype}
\usepackage{enumitem,xurl,needspace}
\usepackage[dvipsnames]{xcolor}
\usepackage{hyperref}
\hypersetup{unicode=true,colorlinks=true,linkcolor=MidnightBlue,urlcolor=MidnightBlue,
 pdftitle={UnsolvedMath Problem 2146},pdfauthor={Source-annotated compilation},bookmarksnumbered=true}
\usepackage{bookmark,tocloft,titlesec,fancyhdr}
\setlength{\cftsecnumwidth}{4.2em}
\cftsetpnumwidth{2.5em}
\cftsetrmarg{4em}
\titleformat{\section}{\Large\bfseries\raggedright}{\thesection}{0.7em}{}
\pagestyle{fancy}\fancyhf{}
\fancyhead[L]{\small\sffamily Additional UnsolvedMath statements}
\fancyhead[R]{\small Original catalogue IDs}
\fancyfoot[C]{\thepage}
\setlength{\parindent}{0pt}
\setlength{\parskip}{5pt plus 1pt}
\setlength{\emergencystretch}{3em}
\setcounter{tocdepth}{1}\setcounter{secnumdepth}{1}
\setlist[itemize]{leftmargin=3em,itemsep=4pt,topsep=4pt,parsep=0pt}
\allowdisplaybreaks
\newcommand{\N}{\mathbb{N}}
\newcommand{\Z}{\mathbb{Z}}
\newcommand{\Q}{\mathbb{Q}}
\newcommand{\R}{\mathbb{R}}
\newcommand{\C}{\mathbb{C}}
\newcommand{\F}{\mathbb{F}}
\newcommand{\A}{\mathbb{A}}
\newcommand{\PP}{\mathbb{P}}
\newcommand{\E}{\mathbb{E}}
\newcommand{\eps}{\varepsilon}
\newcommand{\dd}{\,\mathrm d}
\newcommand{\sref}[2]{\hyperlink{source:#1:#2}{[S#2]}}
\newif\ifshowcatalogue
\showcataloguetrue
\newcommand{\cataloguescope}[1]{\ifshowcatalogue
 \paragraph{Statement recorded in the supplied source.}
 #1\fi}
\begin{document}
% UnsolvedMath ID 2146; source catalogue code EP-501

\setcounter{section}{2145}

\section{Independent Sets for Small-Measure Set Mappings}\label{2146}

{\small\sffamily UnsolvedMath ID 2146 · Combinatorics}

\subsection{Definitions and mathematical statement}

\paragraph{Shared notation.}Unless a section says otherwise, $\N=\{1,2,\ldots\}$,
$\Z$ is the ring of integers, $\Q,\R,\C$ are the rational, real, and complex
fields, $[N]=\{1,\ldots,\lfloor N\rfloor\}$, and $\log$ is the natural
logarithm. For real functions of a parameter tending to infinity,
$f\ll g$ and $f=O(g)$ mean $|f|\leq Cg$ eventually for some fixed $C$;
$f\gg g$ reverses this comparison; $f=o(g)$ means $f/g\to0$;
$f\sim g$ means $f/g\to1$; and $f\asymp g$ means both order bounds.
A subscript on $O$ or $\ll$ specifies allowed constant dependence.
The expression $n^{o(1)}$ in an upper bound means growth smaller than
$n^\epsilon$ for each fixed $\epsilon>0$. Local definitions govern any
exception, finite-field convention, or use of the same symbol for another
object. An asymptotic claim is not automatically an effective algorithm.



Outer measure is Lebesgue outer measure; measurability of $A_x$ is not required in the first assertion. Independence imposes both ordered prohibitions for each distinct pair, since $x\notin A_y$ is required for all distinct $x,y\in X$. \sref{2146}{1} \sref{2146}{2}

\paragraph{Statement recorded in the supplied source.}
For every $x\in\mathbb{R}$ let $A_x\subset \mathbb{R}$ be a bounded set with outer measure $<1$. Must there exist an infinite independent set, that is, some infinite $X\subseteq \mathbb{R}$ such that $x\not\in A_y$ for all $x\neq y\in X$?
If the sets $A_x$ are closed and have measure $<1$, then must there exist an independent set of size $3$? \sref{2146}{1}

\paragraph{Residual task reported by the archive.}

The embedded review (2026-08-17) records: Resolve the unrestricted question in ZFC or establish independence. \sref{2146}{1}

{\small\emph{Transcription note.} Only unambiguous escape/formatting damage and nonmathematical serialization debris were repaired in the displayed source statement.}

\subsection{Short English statement}

Do bounded small-outer-measure forbidden sets assigned to each real point always admit an infinite mutually avoiding set, and what happens for closed forbidden sets?

\subsection{Sources}

{\small

\begin{itemize}

\item[\hypertarget{source:2146:1}{S1}] ULAM.AI, \emph{UnsolvedMath}, supplied \texttt{problems.json}, record 2146 (EP-501), statement and background. The embedded literature-review date is 2026-08-17; that is the archive’s review date, not a fresh certification. Dataset landing page: \url{https://huggingface.co/datasets/ulamai/UnsolvedMath}.

\item[\hypertarget{source:2146:2}{S2}] T. F. Bloom, \emph{Erdős Problems}, Problem 501, \url{https://www.erdosproblems.com/501}. Reference imported from the supplied record; the live problem page was not independently checked for this compilation.

\item[\hypertarget{source:2146:3}{S3}] L. Newelski, J. Pawlikowski and W. Seredyński, Infinite free set for small measure set mappings, Proc. Amer. Math. Soc. 1987. \url{https://doi.org/10.1090/S0002-9939-1987-0870425-6}. Bibliographic information imported from the supplied record.

\item[\hypertarget{source:2146:4}{S4}] Erd\H{o}s, P. and Hajnal, A., Some remarks on set theory. VIII. Michigan Math. J. (1960), 187-191. Bibliographic information imported from the supplied record.

\item[\hypertarget{source:2146:5}{S5}] G\l adysz, S., Bemerkungen "uber die {U}nabh\"{a}ngigkeit der {P}unkte in {B}ezug auf mengenwertige {F}unktionen. Acta Math. Acad. Sci. Hungar. (1962), 199--201. Bibliographic information imported from the supplied record.

\item[\hypertarget{source:2146:6}{S6}] Hechler, S. H., On two problems in combinatorial set theory. Bull. Acad. Polon. Sci. S\'{e}r. Sci. Math. Astronom. Phys. (1972), 429-431. Bibliographic information imported from the supplied record.

\end{itemize}

}

\label{end:2146}
\end{document}
